\scanpage{67}
\mspage{73}
\diagram{\begin{tikzcd}[ampersand replacement=\&]
y_1 \arrow[r,"u"] \& y_0 \arrow[d,"v_0"]\\
y'_1 \arrow[u,"v_1"] \arrow[r,"u'"] \& y'_0
\end{tikzcd}}

un hom de $u$ dans $u'$ est un couple $\left\{\begin{aligned}v_0&;y_0\to y'_0,\\v_1&:y'_1\to y_1\end{aligned}\right.$ tel que $v_0 u v_1=u'$.

Ainsi source et but $s(u)$ et $b(u)$ sont l'une contravariante, l'autre covariante en $u$. On prend pour $q$ le foncteur source, pour $p$ le foncteur but.

Si on applique ceci à $X^\circ$, on trouve
\diagram{\begin{tikzcd}[ampersand replacement=\&]
\& {\unsure{\mathscr D|\mathscr P}(X^\circ)\unsure{-}} \arrow[dl,"p_{X^\circ}"'] \arrow[dr,"q_{X^\circ}"] \& \\
X^\circ \& \& (X^\circ)^\circ=X
\end{tikzcd}}

On aura
\[
\mathrm{Ob}\,\deleted{\unsure{\mathrm{Fl}}}\mathscr D(X^\circ)=\mathrm{Fl}(X^\circ)\simeq\mathrm{Fl}(X)
\]
\diagram{\begin{tikzcd}[ampersand replacement=\&]
y_1 \arrow[r,"u"] \& y_0 \arrow[d,"v_1",very thick] \& \text{dans }X\\
y'_1 \arrow[u,"v_0"] \arrow[r,"u'"] \& y'_0 \& \text{dans }X
\end{tikzcd}}
et $\unsure{=^{\mathrm s}|=}\,$ flèches au sens de $\mathscr D(X)$ et au sens de $\mathscr D(X^\circ)$.

Donc
\[
\left\{\begin{aligned}
\mathscr D(\unsure{X^\circ|X})&=\mathscr D(X^\circ)\\
p_{X^\circ}&=q_X\\
q_{X^\circ}&=p_X
\end{aligned}\right.,
\]
le passage de $X$ à $X^\circ$ revient à échanger $p_X$ et $q_X$.

La fibre de $\mathscr D(X)$ en $x\in\mathrm{Ob}\,X$ est $\simeq(X_{/x})^\circ$, celle en $y^\circ\in\mathrm{Ob}\,X^\circ$ est $\simeq {}_{y\backslash}X$
\[
\left\{\begin{aligned}
p_X^{-1}(\{x\})&\simeq(X_{/x})^\circ\\
\deleted{\unsure{p}}q_X^{-1}(\{y^\circ\})&\simeq {}_{y\backslash}X_{\deleted{\unsure{/y}}}
\end{aligned}\right.
\]

\margin{NB : $p_X$ est $W$-lisse à fibres $W$-asph, donc dans $W^{\unsure{u|w}}$. $q_X$ est $W$-propre à fibres \unsure{$W$-asph} donc dans $W^u$.}
% archive sheet number: 66
